\documentclass[10pt,a4paper,onecolumn]{article}

\usepackage{amsmath,amssymb}
\usepackage{graphicx}
\usepackage[margin=2cm]{geometry}
\usepackage{hyperref}
\usepackage{enumitem}
\usepackage{booktabs}
\usepackage{xcolor}
\usepackage{array}

\hypersetup{colorlinks=true,linkcolor=black,urlcolor=blue,citecolor=blue}
\setlength{\parskip}{0.3ex}

\title{\textbf{Defense Document: Multi-Temporal Deforestation Segmentation \\
With Sentinel-2 and Global Forest Watch Labels} \\
\small Persiapan Sidang / Presentasi Lomba}

\author{Deforest.id Research Team}
\date{}

\begin{document}
\maketitle
\thispagestyle{empty}

\tableofcontents
\newpage

% ======================================================================
\section{Executive Summary (Apa yang Kita Buat?)}
% ======================================================================

Kami membangun \textbf{sistem deteksi deforestasi otomatis} menggunakan
deep learning pada citra satelit Sentinel-2 (resolusi 10~m). Sistem ini
menggunakan \textbf{label dari Global Forest Watch} (GFW, resolusi
30~m) sebagai target training — tanpa label manual.

Kami membandingkan \textbf{tiga arsitektur}:
\begin{enumerate}[nosep]
    \item \textbf{Single-date baseline:} U-Net sederhana, hanya melihat
          citra setelah deforestasi (5 channel). Skor: test IoU 0.122.
    \item \textbf{Stacked multi-temporal:} Pre-event + post-event
          digabung jadi 10 channel input. Skor: test IoU 0.119.
    \item \textbf{Siamese FC-Siam-diff:} Dua encoder identik (weight
          sharing) untuk pre dan post, decoder menghitung perbedaan
          fitur secara eksplisit. Skor: \textbf{test IoU 0.137}.
\end{enumerate}

\textbf{Kesimpulan utama:} Cara kita menggabungkan informasi temporal
\textbf{lebih penting} daripada apakah kita menggunakan informasi
temporal atau tidak. Siamese architecture menang karena explicit change
features.

% ======================================================================
\section{Problem Statement: Kenapa Ini Perlu Dibuat?}
% ======================================================================

\subsection{Masalah Deforestasi di Indonesia}
Indonesia adalah negara dengan \textbf{hutan tropis terluas ketiga}
di dunia ($\sim$94 juta hektar). Setiap tahun, sekitar 300.000 hektar
hutan alam hilang. Dari jumlah ini, sebagian besar adalah
\textbf{small-holder clearing} — pembukaan lahan skala kecil oleh
individu atau komunitas lokal.

\textbf{Mengapa small-holder clearing sulit dideteksi?}
\begin{itemize}[nosep]
    \item Luas lahan yang dibuka: 0.1 -- 0.5 hektar (seukuran
          lapangan basket sampai lapangan sepak bola mini).
    \item Sistem monitoring eksisting (GFW, SIMONTANA) menggunakan
          citra Landsat 30~m. Satu pixel Landsat = 0.09 ha. Jadi
          clearing 0.1 ha hanya muat 1--2 pixel — tidak cukup untuk
          analisis shape atau konfirmasi.
    \item Dilakukan secara tersebar, bukan dalam blok besar.
\end{itemize}

\subsection{Kesenjangan (Research Gap)}
\textbf{Sistem yang ada sekarang:}
\begin{itemize}[nosep]
    \item \textbf{GFW}: 30~m, tahunan. Bagus untuk skala global, tapi
          tidak bisa deteksi sub-hektar.
    \item \textbf{SIMONTANA} (Indonesia): Minimum Mapping Unit 6.25~ha.
          Jelas terlalu besar.
    \item \textbf{Planet Labs} (3~m): Berbayar. Tidak feasible untuk
          monitoring kontinu di Indonesia.
\end{itemize}

\textbf{Solusi yang kami usulkan:}
Gunakan \textbf{Sentinel-2} (10~m, gratis) + \textbf{deep learning}
untuk mendeteksi deforestasi pada skala yang sebelumnya tidak
terjangkau.

\subsection{Pertanyaan Riset}
\begin{enumerate}[nosep]
    \item Bisakah label GFW (30~m) digunakan untuk melatih model
          segmentasi pada Sentinel-2 (10~m)?
    \item Apakah multi-temporal fusion lebih baik daripada
          single-date? Kalau iya, arsitektur fusion mana yang terbaik?
    \item Seberapa besar dampak label noise terhadap performa?
\end{enumerate}

% ======================================================================
\section{Dataset: Dari Mana Datanya, Kenapa?}
% ======================================================================

\subsection{Mengapa Sentinel-2, Bukan Landsat?}

\begin{table}[ht]
\centering
\small
\begin{tabular}{lcc}
\toprule
\textbf{Parameter} & \textbf{Sentinel-2} & \textbf{Landsat 8-9} \\
\midrule
Resolusi spasial (RGB, NIR) & 10~m & 30~m \\
Resolusi temporal & 5 hari & 16 hari \\
Luas minimum detectable & 0.01 ha & 0.09 ha \\
Biaya & Gratis & Gratis \\
Cakupan & Global & Global \\
\bottomrule
\end{tabular}
\end{table}

\textbf{Resolusi 10~m itu penting kenapa?} Deforestasi skala kecil
(0.1 ha) terdiri dari sekitar 100 pixel di Sentinel-2, tapi cuma
$\sim$11 pixel di Landsat. Dengan 100 pixel, model bisa belajar
\textit{tekstur} dan \textit{bentuk} — bukan cuma ada/tidak ada
perubahan.

\textbf{Temporal resolution 5 hari itu penting kenapa?} Indonesia
sering berawan. Dengan revisit 5 hari, kita punya banyak kesempatan
mendapatkan scene bebas awan dalam satu tahun. Landsat 16 hari:
kadang cuma 1--2 scene bersih per tahun.

\subsection{Mengapa 5 Band Input (R, G, B, NIR, NDVI)?}

\textbf{Citra satelit itu bukan foto biasa.} Setiap "band" adalah
lapisan informasi spektral yang berbeda:
\begin{itemize}[nosep]
    \item \textbf{RGB (Blue 490nm, Green 560nm, Red 665nm):} Sama
          seperti mata manusia. Untuk visualisasi dan interpretasi
          visual. Juga menangkap informasi tanah terbuka (soil
          exposure setelah deforestasi).
    \item \textbf{NIR (Near-Infrared, 842nm):} Tidak terlihat oleh
          manusia. Vegetasi sehat memantulkan NIR sangat kuat
          ($\sim$40--50\%). Tanah dan air memantulkan NIR rendah
          ($\sim$10--20\%). Inilah band paling penting: deforestasi
          menyebabkan penurunan NIR yang sangat mencolok.
    \item \textbf{NDVI (Normalized Difference Vegetation Index):}
          Rumusnya $(NIR - Red) / (NIR + Red)$. Hasilnya angka antara
          -1 sampai 1. Hutan lebat: 0.7--0.9. Tanah terbuka: 0.1--0.2.
          Air: negatif. NDVI memberikan informasi vegetasi yang
          \textbf{dinormalisasi} — tidak terpengaruh oleh kecerahan
          absolut citra (bayangan, sudut matahari, dll).
\end{itemize}

\textbf{Mengapa tidak 10 band (semua band Sentinel-2)?}
Sentinel-2 punya 13 band. Band 1 (60m), 2-4 (10m), 5-7 (20m),
8 (10m), 8a (20m), 9-10 (60m), 11-12 (20m). Band 20~m perlu
upsampling ke 10~m. Ini:
\begin{itemize}[nosep]
    \item Menambah ukuran data 2$\times$ tanpa informasi baru (resolusi
          asli sudah 20m).
    \item Menambah noise dari interpolasi.
    \item Uji coba awal dengan 10 band menunjukkan performa setara
          dengan 5 band — jadi tidak ada gunanya.
\end{itemize}

\subsection{Mengapa Cloud-Free Composite, Bukan Single Scene?}

\textbf{Single scene:} Satu gambar pada satu tanggal. Masalah:
\begin{itemize}[nosep]
    \item Awan. Indonesia tropis, awan hampir selalu ada.
    \item Bayangan awan. Bisa mirip deforestasi secara spektral.
    \item Haze (kabut asap). Sering terjadi akibat kebakaran hutan.
\end{itemize}

\textbf{Cloud-free composite:} Ambil semua scene dalam 1 tahun,
masking awan, lalu rata-ratakan median/sentinel yang bersih.
Hasilnya: permukaan bumi yang bersih tanpa awan.
\textbf{Trade-off:} Deforestasi yang terjadi di pertengahan tahun
bisa ter-blur karena tercampur dengan citra sebelum dan sesudah
deforestasi. \textbf{Kami menerima trade-off ini} karena model yang
dilatih dengan data berawan akan belajar mendeteksi awan, bukan
deforestasi.

\subsection{Mengapa 64$\times$64 Pixel (0.64 ha)?}

Ini adalah keputusan yang sangat penting. Mari kita analisis:

\begin{itemize}[nosep]
    \item \textbf{GPU Memory:} RTX 3060 punya 12GB VRAM. Dengan
          64$\times$64, batch size 64 muat. Kalau 128$\times$128,
          batch size maksimal cuma 16 — terlalu kecil untuk training
          stabil.
    \item \textbf{Receptive Field:} U-Net dengan 4 encoder block
          (masing-masing stride 2) punya teoritical receptive field
          140$\times$140. Artinya, output di satu pixel sebenarnya
          "melihat" area 140$\times$140 di input. Ini lebih besar
          dari chip 64$\times$64, jadi konteksnya cukup.
    \item \textbf{Small Clearing:} Deforestasi 0.1 ha = 10$\times$10
          pixel di Sentinel-2. Chip 64$\times$64 muat untuk
          mendeteksi ini dengan konteks sekelilingnya.
    \item \textbf{Jumlah Sampel:} 64$\times$64 menghasilkan lebih
          banyak chip per scene daripada 128$\times$128. Kita
          dapatkan 23,135 training chips.
\end{itemize}

\textbf{Mengapa tidak lebih kecil (32$\times$32)?} Terlalu kecil:
tidak cukup konteks untuk membedakan deforestasi dari non-forest.
Boundary informasi hilang.

\subsection{Mengapa GFW Label?}

Ini adalah \textbf{pertanyaan paling kritis}. Kenapa kita tidak pakai
label yang bagus?

\textbf{Realitas:} Membuat label deforestasi manual untuk ribuan
hektar hutan Indonesia membutuhkan:
\begin{itemize}[nosep]
    \item Tenaga annotator terlatih (bisa membedakan deforestasi dari
          logging, kebakaran, pertanian).
    \item Waktu berbulan-bulan.
    \item Biaya puluhan juta rupiah.
\end{itemize}

\textbf{GFW (Global Forest Watch)} \cite{hansen2013high} menyediakan
\textit{lossyear} layer: peta tahunan kehilangan tutupan pohon dari
tahun 2000--2023. Ini adalah \textbf{satu-satunya dataset deforestasi
open-access yang mencakup Indonesia}.

\textbf{Tapi GFW punya masalah:}
\begin{itemize}[nosep]
    \item Berasal dari Landsat (30~m). Jadi 1 pixel GFW = 9 pixel
          Sentinel-2.
    \item Ini menciptakan \textbf{label noise struktural}: batas
          deforestasi di GFW tidak presisi di resolusi 10~m.
    \item Ada juga \textbf{temporal noise}: GFW melaporkan perubahan
          tahunan, sementara komposit kami spesifik ke 2020 dan 2023.
          Deforestasi yang terjadi 2021 atau 2022 tercatat sebagai
          "perubahan" tapi timing-nya tidak tepat.
\end{itemize}

\textbf{Mengapa tetap pakai GFW?} Karena ini adalah \textit{proof of
concept}. Tujuan kami bukan mencapai akurasi maksimal, tapi
\textbf{membuktikan pipeline end-to-end bisa jalan}. Kalau pipeline
ini berhasil dengan label noise tinggi, maka dengan label yang lebih
baik (manual, resolusi tinggi), hasilnya pasti lebih baik.

\subsection{Dataset Statistics}

\begin{itemize}[nosep]
    \item \textbf{Sumber data:} 13 scene Sentinel-2 mencakup hutan
          lindung di Riau dan Kalimantan Tengah.
    \item \textbf{Total paired chips (pre+post):} 40,484.
    \item \textbf{Train / Val / Test:} 23,135 / 5,822 / 11,527.
    \item \textbf{Global positive pixel ratio:} 0.38\%.
    \item \textbf{Chips dengan zero deforestation:} 74\% (30,000
          chips dari 40,484).
    \item \textbf{Perbandingan train vs test:} Test set punya density
          deforestasi lebih tinggi (0.9\% vs 0.38\%) — ini penting
          nanti untuk interpretasi hasil.
\end{itemize}

\subsection{Pertanyaan Potensial di Sidang}

\begin{itemize}
    \item{\textbf{Q:} \textit{Kenapa train-test split-nya tidak
          random? Kenapa dipisah per scene?}}
    \item{\textbf{A:} Kalau random, chips dari scene yang sama bisa
          masuk train dan test. Model akan "menghafal" karakteristik
          scene itu. Dengan split geografis, kita menguji
          generalization: apakah model bisa mendeteksi deforestasi
          di lokasi yang belum pernah dilihat? Ini lebih realistis.}
\end{itemize}

\begin{itemize}
    \item{\textbf{Q:} \textit{13 scene itu seberapa representatif
          untuk Indonesia?}}
    \item{\textbf{A:} Terbatas. Riau dan Kalteng adalah hot-spot
          deforestasi, tapi tipe hutan di Papua, Kalimantan Utara,
          Sulawesi berbeda. Generalizability masih perlu diuji.} 
\end{itemize}

% ======================================================================
\section{Teori Dasar: Biar Paham Betul}
% ======================================================================

\subsection{Apa Itu Image Segmentation?}
Image segmentation = memberi label ke \textbf{setiap pixel} dalam
gambar. Bukan cuma "apakah gambar ini ada hutan?" tapi "pixel mana
saja yang merupakan hutan?".

Dalam kasus kami: setiap pixel output adalah probabilitas antara 0
(bukan deforestasi) dan 1 (deforestasi). Threshold di 0.5 untuk
dapatkan mask biner.

\subsection{Apa Itu U-Net?}
U-Net \cite{ronneberger2015unet} adalah arsitektur deep learning
untuk segmentasi, bentuknya seperti huruf U:
\begin{itemize}[nosep]
    \item \textbf{Sisi kiri (encoder):} Konvolusi + pooling.
          Mengekstrak fitur dari gambar, dari detail halus (tepi,
          tekstur) ke abstrak (bentuk, objek). Resolusi mengecil
          (64$\times$64 $\rightarrow$ 32$\times$32 $\rightarrow$
          16$\times$16 $\rightarrow$ 8$\times$8 $\rightarrow$
          4$\times$4).
    \item \textbf{Dasar (bottleneck):} Fitur paling abstrak. Model
          "memahami" apa yang ada di gambar secara global.
    \item \textbf{Sisi kanan (decoder):} Upsampling + konvolusi.
          Mengembalikan resolusi ke ukuran asli. Menggabungkan fitur
          abstrak dari bottleneck dengan fitur detail dari encoder
          (via \textbf{skip connections}).
    \item \textbf{Skip connections:} Ini inovasi utama U-Net.
          Decoder tidak hanya bekerja dari bottleneck, tapi juga
          menerima informasi langsung dari encoder di level yang
          sama. Ini mempertahankan detail spasial (posisi, batas).
\end{itemize}

\subsection{Apa Itu ResNet34?}
ResNet \cite{he2016deep} adalah encoder yang menggunakan
\textbf{residual blocks}: alih-alih belajar $F(x)$ langsung, belajar
$F(x) = H(x) - x$ (residual). Keuntungan:
\begin{itemize}[nosep]
    \item Gradient bisa mengalir langsung melalui shortcut connections
          (mencegah vanishing gradient di jaringan dalam).
    \item 34 layer dengan 21.3M parameter.
    \item Tersedia pretrained di ImageNet (14 juta gambar natural).
\end{itemize}

\subsection{Apa Itu Siamese Network?}
Siamese = "kembar". Dua subnet dengan \textbf{weight yang sama persis}
(shared weights). Input berbeda (pre dan post), encoder identik.
Keluaran encoder dibandingkan.

\textbf{Kenapa weight sharing?} Karena kita ingin encoder mengekstrak
fitur yang konsisten. Kalau encoder pre dan post berbeda, model bisa
"curang": satu encoder belajar deteksi hutan, satu lagi belajar
deteksi non-hutan. Dengan weight sharing, encoder \textit{harus}
mengekstrak fitur yang sama dari kedua input — perbedaan hanya muncul
dari inputnya.

\subsection{Apa Itu IoU (Intersection over Union)?}
\textbf{Metric paling penting untuk segmentasi.}

$$IoU = \frac{TP}{TP + FP + FN}$$

\begin{itemize}[nosep]
    \item \textbf{TP (True Positive):} Pixel yang benar-benar
          deforestasi dan diprediksi sebagai deforestasi.
    \item \textbf{FP (False Positive):} Pixel yang bukan deforestasi
          tapi diprediksi deforestasi (false alarm).
    \item \textbf{FN (False Negative):} Pixel deforestasi yang tidak
          terdeteksi (missed detection).
\end{itemize}

IoU = 1 artinya prediksi sempurna. IoU = 0 artinya tidak ada overlap.
Dalam segmentasi satelit, IoU 0.5 sudah dianggap bagus untuk objek
kecil.

\subsection{Apa Itu Dice Loss?}
Dice loss adalah fungsi loss yang digunakan untuk training:
$$DiceLoss = 1 - \frac{2 \times TP + \epsilon}{2 \times TP + FP + FN + \epsilon}$$

Ini adalah \textbf{smoothing} dari IoU. Bedanya: TP dihitung 2$\times$
sehingga lebih toleran terhadap false positives. Smoothing $\epsilon$
mencegah division by zero.

\textbf{Mengapa efektif untuk imbalance?} Karena hanya menghitung
berdasarkan TP, FP, FN — \textbf{tidak dipengaruhi oleh TN (True
Negative)}. Dalam dataset dengan 99.62\% background, TN mendominasi
total pixel. Cross-entropy loss akan puas dengan prediksi
all-background karena akurasi 99.62\%. Dice loss tidak peduli dengan
background yang benar — hanya peduli dengan \textit{overlap} terhadap
kelas positif.

% ======================================================================
\section{Arsitektur: Kenapa Tiga Model?}
% ======================================================================

\subsection{Pertanyaan Fundamental}
Untuk multi-temporal deforestation, ada \textbf{dua pertanyaan}:
\begin{enumerate}[nosep]
    \item Apakah menggunakan informasi pre-event \textit{sama sekali}
          membantu dibandingkan hanya post-event?
    \item Kalau ya, \textbf{bagaimana cara terbaik} menggabungkan
          pre dan post?
\end{enumerate}

Untuk menjawabnya kita butuh 3 model:

\subsection{Model 1: Single-Date (Baseline)}
\begin{itemize}[nosep]
    \item Input: 5 channel (post-event only: R, G, B, NIR, NDVI).
    \item Arsitektur: SimpleUNet (ResNet34-style, 21.3M params).
    \item \textbf{Mengapa ini baseline?} Ini mewakili \textit{status
          quo}: kebanyakan sistem deteksi deforestasi menggunakan
          satu citra saja. Lihat hutan, lihat bukan hutan. Tidak ada
          informasi perubahan.
    \item \textbf{Keterbatasan:} Tidak bisa membedakan "lahan yang
          selalu terbuka" dengan "lahan yang baru dibuka". Kebun
          sawit, ladang kering, tanah kosong — semua terlihat mirip
          dengan deforestasi di satu titik waktu.
\end{itemize}

\subsection{Model 2: Stacked Multi-Temporal}
\begin{itemize}[nosep]
    \item Input: 10 channel (5 pre + 5 post, di-concatenate).
    \item Arsitektur: Sama dengan baseline, layer pertama diadaptasi
          dari 3 channel (pretrained) ke 10 channel.
    \item \textbf{Mengapa ada model ini?} Untuk menguji: apakah
          dengan memberikan informasi pre dan post \textit{tanpa
          struktur khusus}, model bisa belajar perubahan?
    \item \textbf{Masalah:} Model menerima 10 channel, tapi tidak
          tahu mana pre dan mana post. Channel 1-5 pre, 6-10 post.
          Model harus belajar sendiri bahwa perbedaan channel 1-5
          vs 6-10 itu signifikan. Ini tugas tambahan yang tidak perlu.
\end{itemize}

\subsection{Model 3: Siamese FC-Siam-diff}
\begin{itemize}[nosep]
    \item Input: Pre (5 channel) dan Post (5 channel) diproses
          \textbf{terpisah}.
    \item Arsitektur: Dua encoder ResNet34 dengan weight sharing.
          Decoder menerima $[f_{post}, f_{pre}, f_{post} - f_{pre}]$.
    \item \textbf{Inovasi:} Perbedaan fitur dihitung \textbf{secara
          eksplisit} di setiap level decoder. Model tidak perlu
          menebak mana yang berubah — perubahannya sudah disediakan.
    \item \textbf{Mengapa weight sharing?} Agar encoder menghasilkan
          representasi yang konsisten. Perbedaan output encoder
          semata-mata karena perbedaan input (pre vs post), bukan
          karena encoder berbeda.
\end{itemize}

\subsection{Mengapa Bukan Model Lain?}

\begin{itemize}
    \item{\textbf{Vision Transformer (ViT):} Butuh data sangat besar
          ($>$100k images) untuk performa baik. Kita cuma punya 23k
          chips, dan itu pun 74\% tanpa deforestasi.}
    \item{\textbf{DeepLabv3+:} Arsitektur kuat, tapi atrous
          convolution-nya (dilated convolution) boros memori untuk
          chip kecil 64$\times$64. Keuntungan ASPP (Atrous Spatial
          Pyramid Pooling) untuk multi-scale context kurang berguna
          pada chip kecil.}
    \item{\textbf{SegNet:} Lebih efisien memori (tanpa skip
          connections menyimpan banyak data), tapi akurasi boundary
          lebih rendah. Untuk deteksi deforestasi, boundary presisi
          itu penting.}
    \item{\textbf{Attention U-Net:} Menambahkan attention gate di
          skip connections. Ini interessante, tapi menambah
          kompleksitas. Prioritas kami adalah baseline multi-temporal
          dulu.}
\end{itemize}

\subsection{Pertanyaan Potensial di Sidang}

\begin{itemize}
    \item{\textbf{Q:} \textit{Kenapa ResNet34 bukan ResNet50?}}
    \item{\textbf{A:} ResNet50 (25.6M params) tidak muat di memori
          dengan batch 64. Selain itu, untuk chip 64$\times$64,
          representasi ResNet34 sudah cukup. Bottleneck ResNet50
          (1$\times$1, 3$\times$3, 1$\times$1) baru berguna untuk
          gambar lebih besar ($>$224$\times$224).}
\end{itemize}

\begin{itemize}
    \item{\textbf{Q:} \textit{Bukannya weight sharing menghambat
          representasi spesifik?}}
    \item{\textbf{A:} Weight sharing memaksa encoder mengekstrak
          fitur yang sama dari pre dan post. Justru ini yang kita
          mau: perbedaan deteksi harus berasal dari input, bukan dari
          encoder. Analogi: dua foto diambil kamera yang sama —
          perbedaan hasil karena subjeknya berubah, bukan karena
          kameranya berbeda.}
\end{itemize}

\begin{itemize}
    \item{\textbf{Q:} \textit{Ada opsi arsitektur temporal lain?}}
    \item{\textbf{A:} Banyak. ConvLSTM, 3D CNN, Temporal Attention
          Transformer. Tapi untuk 2 waktu (pre/post), Siamese sudah
          optimal. Arsitektur temporal kompleks baru berguna untuk
          time-series multi-waktu ($>$5 timestep).}
\end{itemize}

% ======================================================================
\section{Training: Semua Detail + Alasannya}
% ======================================================================

\subsection{Fungsi Loss: Kenapa Dice?}

\textbf{Pertanyaan inti:} Dataset kita punya 0.38\% pixel positif,
99.62\% background. Loss function yang tidak dirancang untuk imbalance
akan gagal total.

\textit{Kita coba jelaskan dengan analogi:} Bayangkan kita mencari
jarum di tumpukan jerami. Cross-entropy = ukur seberapa benar kita
mengidentifikasi jerami. Dice loss = ukur seberapa baik kita
menemukan jarum.

\textbf{Matematika sederhana:}
\begin{itemize}[nosep]
    \item \textbf{Cross-Entropy:} $CE = -[y \log(p) + (1-y) \log(1-p)]$.
          Untuk setiap pixel, hitung loss-nya. Background (99.62\%)
          mendominasi total loss. Model optimal: prediksi semua 0
          (background), CE rendah, IoU = 0.
    \item \textbf{Dice Loss:} $Dice = 1 - \frac{2TP}{2TP+FP+FN}$.
          Hanya pixel yang diprediksi positif berkontribusi. Tidak
           peduli background sebanyak apapun.
\end{itemize}

\textbf{Uji coba loss lain:}
\begin{itemize}[nosep]
    \item \textbf{Focal Loss:} Modifikasi CE yang menekan loss dari
          easy samples. Hasil: lebih buruk dari Dice. Kemungkinan
          karena hyperparameter $\gamma$ dan $\alpha$ sulit dituning
          untuk extreme imbalance.
    \item \textbf{Tversky Loss:} Variasi Dice dengan bobot berbeda
          untuk FP dan FN. Dengan $\alpha=0.7$ (bobot FN lebih
          besar), recall naik tapi precision turun drastis. Untuk
          aplikasi yang membutuhkan kepercayaan tinggi (alert
          system), precision juga penting.
    \item \textbf{Lov\'asz-Softmax:} Optimasi langsung IoU. Hasil:
          setara dengan Dice tapi training lebih lambat (perlu
          sorting probabilities per batch).
\end{itemize}

\subsection{Optimizer: Kenapa AdamW?}

\textbf{Adam} (Adaptive Moment Estimation):
\begin{itemize}[nosep]
    \item Setiap parameter punya learning rate sendiri berdasarkan
          history gradient-nya. Cocok untuk masalah sparse gradient
          (banyak pixel background, sedikit foreground).
    \item Lebih stabil daripada SGD untuk loss landscape yang kompleks.
\end{itemize}

\textbf{AdamW = Adam + Weight Decay yang benar:}
\begin{itemize}[nosep]
    \item Weight decay = menambahkan $\lambda \cdot w$ ke gradient.
          Efek: parameter yang tidak berguna menuju nol (regularisasi).
    \item Di Adam biasa, weight decay tercampur dengan adaptive
          gradient — tidak efektif. AdamW memisahkan keduanya.
    \item Weight decay $5\times 10^{-4}$: nilai standar untuk
          ResNet. Terlalu besar ($10^{-2}$): underfitting. Terlalu
          kecil ($10^{-5}$): tidak ada efek.
\end{itemize}

\subsection{Learning Rate Scheduling: Kenapa Cosine Annealing?}

\textbf{Learning rate} menentukan seberapa besar langkah update
parameter.

\begin{itemize}[nosep]
    \item \textbf{Awal training:} LR tinggi ($10^{-4}$) = eksplorasi
          cepat. Model bergerak cepat di loss landscape.
    \item \textbf{Akhir training:} LR rendah (menuju 0) = konvergensi
          halus. Model mencari minimum lokal yang baik.
\end{itemize}

\textbf{Cosine Annealing:} LR mengikuti kurva cosinus dari
$10^{-4} \rightarrow 0$ selama 30 epoch. Transisi smooth.

\textbf{Alternatif yang tidak dipilih:}
\begin{itemize}[nosep]
    \item \textbf{Step decay} (turun di epoch 10, 20): Penurunan
          mendadak bisa mengganggu konvergensi.
    \item \textbf{Reduce on plateau} (turun kalau val loss stagnan):
          Butuh patience tuning. Untuk training dengan 30 epoch,
          tidak cukup waktu untuk plateau detection.
    \item \textbf{One-cycle}: Menaikkan LR dulu baru turun. Bagus
          untuk fast training tapi butuh tuning max LR.
\end{itemize}

\subsection{Gradient Accumulation: Kenapa Accum=2?}

\textbf{Masalah:} RTX 3060 hanya muat batch size 64. Tapi batch 64
terlalu kecil untuk training stabil dengan class imbalance.

\textbf{Solusi:} Gradient accumulation. Proses:
\begin{enumerate}[nosep]
    \item Forward pass batch 1 (64 chips), hitung loss, hitung
          gradient. \textbf{Tapi jangan update parameter.}
    \item Forward pass batch 2 (64 chips), hitung loss, akumulasi
          gradient. \textbf{Baru update parameter.}
\end{enumerate}

\textbf{Efek:} Gradient dari 128 chips diakumulasi sebelum update.
Sama seperti batch size 128, tapi tanpa perlu memori 2$\times$.

\textbf{Mengapa accumulation 2, bukan 4 atau 8?} Accum 2 sudah
cukup untuk mencegah collapse. Accum 4 (efektif 256) tidak memberikan
peningkatan signifikan. Accum 8 tidak muat di memori (perlu nyimpen
activation untuk 512 chips).

\subsection{Oversampling: Kenapa 10$\times$?}

\textbf{Masalah:} 74\% chips tanpa deforestasi. Sampling acak:
rata-rata hanya 17 dari 64 chips per batch mengandung deforestasi.

\textbf{Oversampling:} Setiap chip positif di-sampling 10$\times$
lebih sering. Cara: dataset dibagi jadi positif (26\%) dan negatif
(74\%). Dari positif, sampling dengan replacement 10$\times$ ukuran
asli. Gabung dengan negatif, shuffle.

\textbf{Efek:} Proporsi chips positif dalam batch $\sim$50\%. Model
cukup sering melihat deforestasi untuk belajar.

\textbf{Mengapa 10, bukan 5 atau 20?}
\begin{itemize}[nosep]
    \item \textbf{5$\times$:} Cuma 30\% batch positif. Masih terlalu
          sedikit — model tetap bias ke background.
    \item \textbf{20$\times$:} Dominasi chips positif yang sama
          (di-sampling berulang). Model overfit ke beberapa contoh
          spesifik.
    \item \textbf{10$\times$:} Proporsi $\sim$50:50. Ideal.
\end{itemize}

\subsection{Augmentasi: Kenapa Minimal?}

Data augmentation = memvariasikan data training secara artifisial.
Tujuan: meningkatkan generalisasi.

\textbf{Kami hanya menggunakan flip horizontal + vertikal.}
Mengapa?
\begin{itemize}[nosep]
    \item \textbf{Spectral consistency sangat penting.} Deforestasi
          dideteksi dari perubahan spektral (kandungan NIR turun,
          NDVI turun). Rotasi mengubah sudut datang cahaya, brightness
          mengubah intensitas, Gaussian noise mengubah distribusi
          spektral. Semua ini merusak sinyal deforestasi.
    \item \textbf{Definisi "berubah" itu absolut.} Kalau kita
          mengubah kecerahan pre dan post secara independen, selisih
          NDVI sebelum-sesudah juga berubah — padahal di lapangan
          tidak ada perubahan kecerahan.
    \item \textbf{Eksperimen validasi:} Heavy augmentation (rotation
          $\pm 30^\circ$, brightness $\pm 10\%$, noise $\sigma=0.01$)
          menurunkan val IoU dari 0.044 ke 0.031. Penurunan 30\%.
\end{itemize}

\textbf{Flip aman:} Deforestasi tidak memiliki orientasi — hutan
yang dibuka dari utara atau selatan tetap deforestasi. Flip tidak
mengubah nilai spektral.

\subsection{AMP (Mixed Precision): Kenapa?}

\textbf{FP32 vs FP16:}
\begin{itemize}[nosep]
    \item FP32: 32-bit floating point. Presisi tinggi, memori besar.
    \item FP16: 16-bit floating point. Presisi setengah, memori
          setengah, throughput 2$\times$ di Tensor Cores.
\end{itemize}

\textbf{AMP} = Automatic Mixed Precision. Bagian kritis (loss
computation, gradient scaling) tetap FP32. Bagian komputasi berat
(konvolusi) di FP16.

\textbf{Keuntungan untuk RTX 3060:}
\begin{itemize}[nosep]
    \item Tensor Cores (generasi Ampere) memberikan 2-3$\times$
          throughput di FP16 vs FP32.
    \item Memory usage turun 30\% — memungkinkan batch 64 atau
          model dengan lebih banyak channel.
    \item Loss scaling mencegah underflow gradient (gradient terlalu
          kecil untuk FP16).
\end{itemize}

\subsection{Gradient Clipping: Kenapa 1.0?}

\textbf{Problem:} Dengan Dice loss dan imbalance, kadang ada batch
yang loss-nya sangat tinggi (misal batch berisi deforestasi yang
sulit). Gradient-nya bisa sangat besar, menyebabkan parameter update
yang ekstrem — merusak representasi.

\textbf{Solusi:} Clip gradient norm ke 1.0. Artinya: kalau panjang
vektor gradient $>$ 1.0, skala ulang ke 1.0. Informasi arah gradient
tetap, besarnya dibatasi.

\textbf{Kenapa 1.0?} Terlalu kecil (0.1) menghambat learning.
Terlalu besar (10.0) tidak efektif. 1.0 adalah nilai default yang
umum.

\subsection{RAM Caching: Kenapa?}

\textbf{Tanpa cache:} Setiap epoch, 23,135 chips dibaca dari SSD.
Setiap chip: load file NPZ (10 bands $\times$ 64$\times$64), parse,
stack ke tensor. Total: $\sim$210 detik/epoch. 30 epoch = 105 menit
hanya I/O.

\textbf{Dengan cache:} Epoch 1: load semua ke RAM (24 GB). Epoch
2-30: baca dari RAM. $\sim$27 detik/epoch. 30 epoch = 16 menit.

\textbf{Mengapa feasible?} 23,135 chips $\times$ 10 bands $\times$
64 $\times$ 64 $\times$ 4 bytes = $\sim$3.8 GB. Komputer punya 32 GB
RAM, 24 GB free setelah OS. Muat.

% ======================================================================
\section{Inference \& Evaluation}
% ======================================================================

\subsection{Test-Time Augmentation (TTA)}

\textbf{Ide:} Saat inference, jangan hanya jalankan sekali. Jalankan
3 kali (original, flip horizontal, flip vertical), rata-rata
probabilitasnya.

\textbf{Kenapa efektif?} Ensemble 3 model (dengan parameter sama
tapi input berbeda). Setiap varian flip punya "kesalahan" berbeda.
Rata-rata menghaluskan kesalahan.

\textbf{Hasil:}
\begin{itemize}[nosep]
    \item Single-date: IoU 0.119 $\rightarrow$ 0.122 (+2.5\%).
    \item Siamese: IoU 0.132 $\rightarrow$ 0.137 (+3.8\%).
\end{itemize}

\subsection{Global vs Per-Sample Metrics}

Ini \textbf{sangat penting} untuk dipahami. Dua cara menghitung IoU:

\textbf{Per-sample mean:}
\begin{enumerate}[nosep]
    \item Hitung IoU tiap chip: $IoU_1, IoU_2, ..., IoU_n$.
    \item Rata-rata: $\frac{1}{n}\sum IoU_i$.
    \item \textbf{Masalah:} 74\% chips IoU=0 (tanpa deforestasi).
          Mean IoU $\sim$0.013. Angka ini sangat rendah dan tidak
          informatif.
\end{enumerate}

\textbf{Global (summed):}
\begin{enumerate}[nosep]
    \item Jumlahkan semua TP, FP, FN dari semua chips.
    \item $IoU_{global} = \frac{\sum TP}{\sum TP + \sum FP + \sum FN}$.
    \item \textbf{Keuntungan:} Chips tanpa deforestasi tidak
          berkontribusi (TP=0, FP=0, FN=0). Hanya chips dengan
          deforestasi yang memengaruhi metrik.
\end{enumerate}

\textbf{Kapan pakai yang mana?} Per-sample: kalau kita peduli
per-chip (setiap chip harus benar). Global: kalau kita peduli total
area (seberapa luas deforestasi terdeteksi). Untuk aplikasi
monitoring, global lebih relevan.

\subsection{Mengapa Threshold 0.5?}

Dice loss secara implisit mengoptimalkan threshold 0.5. Model
dilatih untuk output probabilitas, dan dice loss optimal saat
threshold 0.5. Kalau kita ubah threshold (misal 0.3), precision
turun, recall naik — trade-off. 0.5 adalah titik natural.

\subsection{Pertanyaan Potensial}

\begin{itemize}
    \item{\textbf{Q:} \textit{Kenapa validation IoU lebih rendah dari
          test IoU?}}
    \item{\textbf{A:} Distribusi deforestasi tidak seragam. Test set
          memiliki proporsi deforestasi lebih tinggi (0.9\%) daripada
          validation set (0.38\%). Karena metrik global dihitung dari
          total TP/FP/FN, semakin banyak deforestasi sebenarnya yang
          ada, semakin tinggi IoU-nya (selama model tidak
          all-background). Ini menunjukkan bahwa model bekerja lebih
          baik di area dengan deforestasi lebih padat.}
\end{itemize}

\begin{itemize}
    \item{\textbf{Q:} \textit{Apakah metrik global tidak menyesatkan?
          Kalau satu chip mendominasi, metrik bisa bias.}}
    \item{\textbf{A:} Valid. Itu sebabnya kami juga melaporkan
          per-sample mean. Tapi untuk tujuan deteksi area, chip besar
          dengan banyak deforestasi harusnya berbobot lebih —
          karena area deforestasi yang lebih luas lebih penting untuk
          dideteksi. Global metrics mencerminkan ini.}
\end{itemize}

% ======================================================================
\section{Hasil: Interpretasi Lengkap}
% ======================================================================

\subsection{Tabel Hasil}

\begin{table}[ht]
\centering
\small
\begin{tabular}{lcccc}
\toprule
\textbf{Model} & \textbf{IoU} & \textbf{Precision} & \textbf{Recall} & \textbf{F1} \\
\midrule
Single-date (post only) & 0.122 & 0.198 & 0.239 & 0.217 \\
Multi-temporal (stacked) & 0.119 & 0.154 & 0.344 & 0.212 \\
\textbf{Siamese (ours)} & \textbf{0.137} & 0.196 & 0.311 & \textbf{0.240} \\
\bottomrule
\end{tabular}
\end{table}

\subsection{Analisis Per Model}

\textbf{Single-date (post only):}
\begin{itemize}[nosep]
    \item Precision tertinggi (0.198). Artinya: dari semua prediksi
          positif, 19.8\% benar.
    \item Recall terendah (0.239). Hanya 23.9\% deforestasi
          sebenarnya yang terdeteksi.
    \item \textbf{Interpretasi:} Model ini konservatif. Hanya
          memprediksi positif kalau sangat yakin. Akibatnya, banyak
          deforestasi terlewat.
    \item \textbf{Kenapa?} Karena tidak punya informasi "sebelumnya",
          model harus menebak dari post-event saja. Area yang
          spektralnya mirip deforestasi (lahan terbuka, pertanian)
          tidak bisa dibedakan.
\end{itemize}

\textbf{Stacked multi-temporal:}
\begin{itemize}[nosep]
    \item Recall tertinggi (0.344). Mendeteksi 34.4\% deforestasi.
    \item Precision terendah (0.154). Banyak false positive.
    \item \textbf{Interpretasi:} Model ini agresif. Concatenation
          10 channel memberi informasi pre+post, tapi model tidak
          bisa memanfaatkannya dengan baik. Akibatnya: perubahan
          spektral \textit{apapun} dianggap deforestasi — termasuk
          perubahan musiman (pertanian, air).
    \item \textbf{Kenapa stacking kalah?} Informasi "mana channel pre
          dan post" hanya dari urutan. Model harus belajar relasi
          antar channel 1-5 vs 6-10. Ini tugas tambahan yang tidak
          perlu. Siamese memberikan relasi ini secara eksplisit.
\end{itemize}

\textbf{Siamese FC-Siam-diff:}
\begin{itemize}[nosep]
    \item IoU tertinggi (0.137), F1 tertinggi (0.240).
    \item Precision (0.196) hampir menyamai single-date.
    \item Recall (0.311) mendekati stacked.
    \item \textbf{Interpretasi:} \textit{Best of both worlds}.
          Explicit difference features memungkinkan model mendeteksi
          perubahan dengan presisi tinggi (seperti single-date) dan
          recall baik (seperti multi-temporal).
    \item \textbf{Kenapa siamese menang?} Karena model tidak perlu
          menebak:
          \begin{enumerate}[nosep]
              \item Encoder weight sharing: representasi konsisten.
              \item Decoder menerima $f_{post} - f_{pre}$: perubahan
                    langsung diberikan, bukan disembunyikan.
              \item Skip connections dengan difference: perubahan di
                    berbagai skala (dari detail tepi hingga fitur
                    global).
          \end{enumerate}
\end{itemize}

\subsection{Mengapa IoU Hanya 0.137?}

\textbf{Pertanyaan paling berat.} Jawaban jujur:

\begin{enumerate}[nosep]
    \item \textbf{Label noise (faktor dominan):} GFW label 30~m,
          citra 10~m. Satu pixel GFW = 9 pixel S2. Misal: deforestasi
          0.5 ha di GFW — area 50$\times$50m — direpresentasikan
          sebagai $\sim$3 pixel Landsat (masing-masing 30$\times$30m).
          Di S2 (10m), ini area 50$\times$50 pixel. Tapi labelnya
          hanya 3 pixel. Model bingung: prediksi besar, label kecil.
          Ini dihitung sebagai FP + FN, menurunkan IoU.
    \item \textbf{Class imbalance:} 0.38\% positif. Dari 1000 pixel,
          hanya 4 yang deforestasi. Model bisa prediksi salah 2 dan
          masih akurasi 99.8\%. Tapi IoU-nya anjlok.
    \item \textbf{74\% chips tanpa deforestasi:} Ini masalah untuk
          training (pengaruh ke loss) dan untuk evaluasi (pengaruh ke
          per-sample metrics).
    \item \textbf{Chip size 64$\times$64:} Terlalu kecil untuk konteks
          landscape. Model tidak bisa membedakan "ini tepi hutan"
          dari "ini deforestasi baru".
    \item \textbf{Cloud-free composite:} Deforestasi yang terjadi di
          pertengahan tahun ter-blur dalam komposit.
\end{enumerate}

\subsection{Perbandingan dengan ICP Paper (NDVI pseudo-label)}

ICP paper (sebelumnya) mencapai IoU 0.73. \textbf{Mengapa jauh
berbeda?}
\begin{itemize}[nosep]
    \item Label ICP = NDVI differencing: $\text{NDVI}_{post} -
          \text{NDVI}_{pre}$. Threshold untuk menentukan deforestasi.
          Ini adalah \textbf{metode hand-crafted} — label dibuat dari
          data yang sama dengan input. Model U-Net belajar meniru
          decision boundary yang sudah ada. Wajar IoU tinggi.
    \item GFW label = \textbf{independent source} dari Landsat.
          Model harus belajar korelasi antara citra S2 10~m dan label
          L8 30~m. Ini tugas yang jauh lebih sulit.
    \item ICP IoU 0.73 = validasi internal (NDVI vs NDVI). GFW IoU
          0.137 = validasi eksternal (S2 vs Landsat). Dua hal berbeda.
\end{itemize}

\subsection{Overfitting: Analisis}

\textbf{Gejala:} Val IoU peak di epoch 4-6, lalu turun. Training
loss terus turun (model menghafal data training). Validation loss
naik atau plateau.

\textbf{Penyebab:}
\begin{itemize}[nosep]
    \item \textbf{Rasio parameter:sampel:} 21.3M parameter, 23k
          training chips (efektif hanya $\sim$6k yang positif).
          Rasio 3,550:1. Idealnya $<$ 100:1.
    \item \textbf{Dataset homogen:} 13 scene dari 2 provinsi.
          Variasi terbatas — model cepat "bosan" dan mulai menghafal.
    \item \textbf{Regularisasi terbatas:} Weight decay 5e-4 saja.
          Dropout akan membantu, tapi menambah waktu training dan
          bisa menurunkan peak performance.
\end{itemize}

\textbf{Solusi potensial (untuk sidang/future work):}
\begin{itemize}[nosep]
    \item \textbf{More data:} Tambah scene dari provinsi lain.
    \item \textbf{Label smoothing:} Mengurangi confidence model
          terhadap label noise.
    \item \textbf{Dropout:} Di decoder (biasanya 0.2-0.5).
    \item \textbf{Early stopping:} Berhenti di epoch 5-6 (peak val
          IoU). Tidak perlu 30 epoch penuh.
\end{itemize}

% ======================================================================
\section{Pipeline: Gambaran Sistem End-to-End}
% ======================================================================

\begin{enumerate}[nosep]
    \item \textbf{Data Export (GEE):} Google Earth Engine.
          Export Sentinel-2 composite (2020, 2023) + GFW lossyear
          layer sebagai GeoTIFF.
    \item \textbf{Chip Generation:} Tile GeoTIFF ke 64$\times$64
          chips. Pair pre/post. Simpan sebagai NPZ (5 bands + mask).
    \item \textbf{Training:} Load chips ke RAM. Oversampling.
          DataLoader dengan augmentasi flip. Model forward/backward
          dengan AMP + gradient accumulation. Log metrics ke CSV.
    \item \textbf{Inference:} Load chip test. TTA (3 flips).
          Threshold 0.5. Hitung metrics (global + per-sample).
          Simpan visualisasi (6-panel plot).
    \item \textbf{Analisis:} Bandingkan 3 model. Plot kurva training.
          Analisis error.
\end{enumerate}

\textbf{Hardware:}
\begin{itemize}[nosep]
    \item GPU: NVIDIA RTX 3060 (12GB VRAM).
    \item CPU: Intel/AMD (cukup untuk data loading).
    \item RAM: 32 GB (24 GB untuk cache).
    \item Storage: SSD (I/O bottleneck tanpa cache).
    \item OS: Windows 11 + WSL2 (Ubuntu).
\end{itemize}

\textbf{Software:}
\begin{itemize}[nosep]
    \item PyTorch 2.7, CUDA 12.x.
    \item AMP (torch.cuda.amp).
    \item Rich (progress bar TUI).
    \item tmux (survive SSH disconnect).
\end{itemize}

\textbf{Performance:}
\begin{itemize}[nosep]
    \item Training 30 epochs: $\sim$15 menit (dengan cache).
    \item Inference 11,527 chips + TTA: $\sim$12 menit.
    \item GPU memory usage: $\sim$8 GB (peak).
\end{itemize}

% ======================================================================
\section{Pertanyaan Sulit yang Mungkin Muncul}
% ======================================================================

\subsection{Tentang Relevansi \ dan Kontribusi}

\begin{itemize}
    \item{\textbf{Q:} \textit{IoU 0.137 — ini sangat rendah. Apa
          gunanya?}}
    \item{\textbf{A:} Ini \textit{proof of concept}. Tujuan kami
          bukan mencapai akurasi operasional, tapi membuktikan
          pipeline end-to-end bisa berjalan dengan label yang ada.
          IoU 0.137 dengan GFW label noise adalah baseline. Dengan
          label manual (yang merupakan langkah logis berikutnya),
          IoU diperkirakan meningkat 3-5$\times$. Nilai dari proyek
          ini adalah: \textbf{(1)} pipeline dari GEE ke visualisasi,
          \textbf{(2)} perbandingan arsitektur, \textbf{(3)} identifikasi
          bottleneck (label noise adalah masalah utama).}
\end{itemize}

\begin{itemize}
    \item{\textbf{Q:} \textit{Apa bedanya dengan pekerjaan yang sudah
          ada?}}
    \item{\textbf{A:} Kebanyakan penelitian deforestasi menggunakan
          label NDVI pseudo-label (seperti ICP paper kami sebelumnya)
          atau label manual. Yang \textbf{unik} dari proyek kami:
          menggunakan GFW label — yang available secara global — untuk
          melatih model 10~m. Ini menunjukkan bahwa GFW, meskipun
          tidak sempurna, bisa menjadi \textit{weak supervision} untuk
          deteksi resolusi lebih tinggi. Kontribusi: \textbf{metode}
          dan \textbf{pipeline}, bukan absolute accuracy.}
\end{itemize}

\begin{itemize}
    \item{\textbf{Q:} \textit{Kenapa tidak langsung pake NDVI
          differencing aja? Itu IoU-nya 0.73.}}
    \item{\textbf{A:} NDVI differencing adalah metode \textit{rules-
          based}: threshold manual yang mungkin tidak generalizable
          ke lokasi lain. Deep learning bisa belajar pola yang lebih
          kompleks dari data multi-band, bukan hanya NDVI. Selain
          itu, proyek ini adalah eksperimen: \textit{can we leverage
          existing GFW labels?} Jawabannya: yes, dengan batasan
          tertentu.}
\end{itemize}

\subsection{Tentang Metode}

\begin{itemize}
    \item{\textbf{Q:} \textit{Mengapa tidak pakai data time-series
          (multi-temporal $>$ 2 waktu)?}}
    \item{\textbf{A:} Terbatas oleh GFW label — hanya tahunan. Kalau
          kita punya 6 time points (2018-2023), labelnya cuma satu:
          "deforestasi terjadi di tahun X". Tidak ada label per-waktu.
          Model tidak bisa belajar temporal dynamics tanpa temporal
          labels. Arsitektur 2-waktu (pre/post) adalah yang paling
          sesuai dengan data yang ada.}
\end{itemize}

\begin{itemize}
    \item{\textbf{Q:} \textit{Bukankah lebih baik pakai unsupervised
          change detection?}}
    \item{\textbf{A:} Unsupervised (seperti PCA, MAD, IR-MAD) tidak
          memerlukan label. Tapi outputnya anomaly score — bukan
          "deforestasi" secara semantik. Tidak bisa membedakan
          deforestasi dari perubahan musiman, banjir, atau kebakaran.
          Supervised learning memberikan \textit{semantic
          understanding}: model belajar \textit{apa itu} deforestasi,
          bukan hanya \textit{apa yang berubah}.}
\end{itemize}

\begin{itemize}
    \item{\textbf{Q:} \textit{Ukuran chip 64$\times$64 itu terlalu
          kecil. Kenapa tidak 256$\times$256?}}
    \item{\textbf{A:} Terbatas GPU. Tapi yang lebih penting:
          eksperimen dengan chip lebih besar membutuhkan re-export
          data dari GEE + re-training + re-inference — siklus yang
          panjang. Dengan 64$\times$64, siklus eksperimen cepat
          ($\sim$15 menit per training). Ini memungkinkan hyperparameter
          search yang ekstensif (9 konfigurasi). Setelah arsitektur
          final dipilih, scaling ke chip lebih besar adalah engineering
          work, bukan research.}
\end{itemize}

\subsection{Tentang Evaluasi}

\begin{itemize}
    \item{\textbf{Q:} \textit{Kenapa tidak pake Kappa atau AUC?}}
    \item{\textbf{A:} Kappa: sensitive terhadap prevalence (proporsi
          kelas). Dengan 0.38\% prevalence, Kappa bisa tinggi meski
          IoU rendah. AUC (ROC): untuk binary classification dengan
          threshold. Segmentasi adalah pixel-level, IoU adalah standar
          de facto untuk segmentation task. Precision/Recall/F1
          melengkapi IoU untuk interpretasi yang lebih nuanced.}
\end{itemize}

\begin{itemize}
    \item{\textbf{Q:} \textit{Apakah hasil ini reproducible?}}
    \item{\textbf{A:} Ya. Semua kode ada di GitHub (public). Semua
          hyperparameter dicatat. Random seed di-set. Satu-satunya
          yang tidak fully reproducible: cloud-free composite dari
          GEE (algoritma median bisa sedikit berbeda tergantung
          versi). Tapi secara umum, hasil harus konsisten dalam
          $\pm 0.01$ IoU.}
\end{itemize}

\subsection{Tentang Aplikasi}

\begin{itemize}
    \item{\textbf{Q:} \textit{Apa rencana implementasi ke depannya?}}
    \item{\textbf{A:} (1) Perbaiki label: manual annotation atau
          Planet imagery (3m) untuk label lebih presisi. (2) Scale:
          256$\times$256 chips, model lebih besar. (3) Operationalize:
          deploy sebagai API yang menerima citra baru dan menghasilkan
          alert. (4) Integrasi SAR Sentinel-1 untuk kemampuan
          all-weather.}
\end{itemize}

\begin{itemize}
    \item{\textbf{Q:} \textit{Apakah model bisa dipakai di daerah
          lain?}}
    \item{\textbf{A:} Belum diuji. Model dilatih di Riau dan Kalteng.
          Generalisasi ke Papua (tipe hutan berbeda), Sumatera Barat
          (topografi berbeda), atau Jawa (fragmentasi tinggi) perlu
          diuji. Tapi karena arsitekturnya general-purpose, kami
          optimis dengan fine-tuning atau tambahan data dari daerah
          target.}
\end{itemize}

% ======================================================================
\section{Limitations: Jawaban Jujur}
% ======================================================================

\begin{enumerate}[nosep]
    \item \textbf{Label noise:} Ini batasan \textbf{fundamental}.
          Selama kita pakai GFW label 30~m, akurasi maksimal terbatas.
          Model tidak bisa lebih presisi dari labelnya.
    \item \textbf{Chip size:} 64$\times$64 membatasi konteks spasial.
          Tapi trade-off dengan GPU memory.
    \item \textbf{Temporal coverage:} Dua waktu (pre/post) dengan
          gap 3 tahun. Deforestasi di 2021 dan 2022 tercampur.
          Idealnya: tahunan atau sub-tahunan.
    \item \textbf{Geographic scope:} Hanya 13 scene dari 2 provinsi.
          Belum mewakili keragaman hutan Indonesia.
    \item \textbf{False positives:} Perubahan non-deforestasi
          (pertanian, badan air musiman, kebakaran) masih menghasilkan
          false positives.
    \item \textbf{Cloud dependence:} Sentinel-2 optical — tidak bisa
          menembus awan. Di musim hujan, bisa berbulan-bulan tanpa
          citra bersih.
\end{enumerate}

% ======================================================================
\section{Future Work: Arah Pengembangan}
% ======================================================================

\begin{enumerate}[nosep]
    \item \textbf{Better labels:} Ini prioritas nomor satu. 500-1000
          chips dengan label manual di resolusi 10~m akan memberikan
          lompatan kualitas. Bisa juga pakai Planet imagery (3~m)
          untuk annotasi.
    \item \textbf{Larger chips:} 128$\times$128 atau 256$\times$256
          dengan gradient accumulation lebih besar. Memberikan konteks
          landscape yang lebih baik.
    \item \textbf{SAR integration:} Sentinel-1 (radar) + Sentinel-2
          (optical). Radar tidak terpengaruh awan. Kombinasi optical
          + radar memberikan coverage temporal yang jauh lebih baik.
    \item \textbf{Multi-temporal time-series:} 5+ time points dengan
          temporal attention atau ConvLSTM. Bisa mendeteksi \textit{
          kapan} deforestasi terjadi, bukan hanya \textit{apakah}.
    \item \textbf{Post-processing:} Connected-component filtering
          (buang prediksi $<$ 0.01 ha), CRF untuk boundary smoothing.
    \item \textbf{Deployment:} API endpoint + dashboard monitoring.
          Integrasi dengan Google Earth Engine untuk alert
          near-real-time.
    \item \textbf{Uncertainty estimation:} Monte Carlo Dropout atau
          Deep Ensemble untuk memberikan confidence score pada setiap
          prediksi.
\end{enumerate}

% ======================================================================
\section{Quick Reference Cards}
% ======================================================================

\subsection*{Arsitektur Comparison}

\begin{table}[ht]
\centering
\small
\begin{tabular}{lccc}
\toprule
\textbf{Aspek} & \textbf{Single-date} & \textbf{Stacked} & \textbf{Siamese} \\
\midrule
Input channels & 5 (post) & 10 (pre+post) & 5+5 (terpisah) \\
Parameters & 21.3M & 21.3M & 25.9M \\
Fusion method & N/A & Concatenation & Explicit diff \\
Pretrained & Adapted & Adapted & Shared \\
Weight sharing & N/A & N/A & Yes \\
Best IoU & 0.122 & 0.119 & \textbf{0.137} \\
Key strength & Precision & Recall & Balance \\
Key weakness & Low recall & Low precision & Complexity \\
\bottomrule
\end{tabular}
\end{table}

\subsection*{Hyperparameter Summary}

\begin{table}[ht]
\centering
\small
\begin{tabular}{lll}
\toprule
\textbf{Hyperparameter} & \textbf{Nilai} & \textbf{Alasan} \\
\midrule
Loss function & Dice (symmetric) & 0.38\% imbalance \\
Optimizer & AdamW & Adaptive + weight decay \\
Learning rate & $10^{-4}$ & Default Adam \\
Weight decay & $5 \times 10^{-4}$ & Regularisasi ringan \\
Scheduler & Cosine annealing & Smooth decay \\
Batch size & 64 (eff. 128) & GPU memory \\
Accumulation & 2 & Simulasi batch 128 \\
Oversampling & 10$\times$ & 74\% zero-deforest chips \\
Augmentation & H/V flip only & Spectral consistency \\
Precision & AMP (FP16) & Tensor Cores speed \\
Gradient clipping & 1.0 & Prevent gradient spikes \\
Epochs & 30 & Overfitting after 6 \\
\bottomrule
\end{tabular}
\end{table}

\subsection*{Defense Script — 60 Second Answers}

\begin{itemize}[nosep]
    \item{\textbf{Apa yang kamu buat?} Sistem deteksi deforestasi
          otomatis dengan deep learning pada Sentinel-2 10~m,
          menggunakan label GFW 30~m sebagai training target.}
    \item{\textbf{Kenapa tiga model?} Untuk menguji apakah multi-
          temporal fusion membantu, dan arsitektur mana yang terbaik.}
    \item{\textbf{Kenapa IoU 0.137 rendah?} Label noise GFW (30m
          vs 10m), extreme class imbalance (0.38\% positif), dan
          74\% chips tanpa deforestasi.}
    \item{\textbf{Mengapa siamese menang?} Karena explicit difference
          features memberikan sinyal perubahan yang lebih bersih
          daripada concatenation.}
    \item{\textbf{Apa kontribusi utama?} Pipeline end-to-end dari
          GEE ke trained model + perbandingan arsitektur + identifikasi
          label noise sebagai bottleneck utama.}
    \item{\textbf{Langkah selanjutnya?} Label manual, chip lebih
          besar, SAR integration.}
\end{itemize}

% ======================================================================
\bibliographystyle{unsrt}
\bibliography{icp-refs}
% ======================================================================

\end{document}
